% !TEX root = ../SANER2027-main-labeling.tex
% Introduction
% Claims to support: headline claim + RQs + contributions
% SANER rewrite (2026-09-24): every paragraph is new text (blue). Numbers come from
% Section IV and paper/tables/{triangulation_all_cells,method_comparison_ci}.csv.

\section{Introduction}
\label{sec:01_intro}
Users and contributors of software projects report bugs, request features, and ask questions through issue trackers such as GitHub Issues~\cite{colavito2024leveraging}. Labeling each issue report with its type speeds up its resolution~\cite{liao2018exploring} and supports issue prioritization~\cite{catolino2019not}, issue assignment~\cite{wu2022spatial}, and developer onboarding~\cite{santos2023tag, aracena2025applying}. Yet labels are rarely used~\cite{kallis2019ticket} and often wrong~\cite{antoniol2008bug, herzig2013s}, and manually assigning them is laborious~\cite{fan2017road}. To address these problems, many studies propose methods for automated issue report classification (\irc) and typically evaluate them on the labels \emph{bug}, \emph{feature}, and \emph{question}~\cite{kallis2019ticket, izadi2022catiss, kallis2024nlbse, colavito2025benchmarking, heo2025study}.

Most \irc\ methods learn from labeled issues. Early work trained classic machine-learning classifiers~\cite{antoniol2008bug, kallis2019ticket, pandey2017automated}, later work fine-tuned pre-trained transformer models~\cite{izadi2022catiss, trautsch2022predicting, colavito2022issue, wang2021well}, and recent work applies large language models (LLMs)~\cite{colavito2025benchmarking, aracena2025applying, heo2025study}. Among LLM-based methods, fine-tuning on labeled issues gives the best reported results~\cite{heo2025study, aracena2025applying, aracena2024applyinglargelanguagemodels}, with Low-Rank Adaptation (LoRA~\cite{hu2022lora}) for open models. Fine-tuning, however, adds a training run that must be repeated to learn from newly labeled issues or to adopt a newer model~\cite{fan2017road, kallis2019ticket}.

LLMs can also use labeled issues without training. Through in-context learning, a model infers a task from a few labeled examples in its prompt~\cite{brown2020language}, and choosing the labeled examples most similar to the query issue can outperform choosing them at random~\cite{liu2022makes, yu2023retrieval}. With such retrieval-augmented few-shot prompting (\rag), a newly labeled issue only needs to be indexed, and the LLM can be replaced without retraining. For future work, Heo and Lee~\cite{heo2025study} suggest few-shot prompting, and Aracena et~al.~\cite{aracena2025applying} suggest retrieval augmentation to reduce the need for extensive fine-tuning and to help with ambiguous labels such as question. To our knowledge, however, retrieval-augmented few-shot prompting has not been systematically evaluated for \irc\ or compared with LoRA fine-tuning.

This paper does both. In an empirical study, we use four Qwen2.5-Instruct models of 3B to 32B parameters~\cite{yang2024qwen25} and Heo and Lee's balanced benchmark of issue reports from eleven projects~\cite{heo2025study}, each with 300 labeled and 300 test issues. Following Heo and Lee, the retrieval index and the fine-tuning data come either from the target project's 300 labeled issues alone (the project-specific setting) or from all 3{,}300 labeled issues of the eleven projects (the project-agnostic setting).

The retrieved examples are the labeled issues most similar to the query, so we first test how well their labels alone predict the query's label. We label the query by a similarity-weighted vote over these neighbors, without an LLM (\knn). Using only the target project's labeled issues, the vote reaches 59.5\% macro~$F_1$, so the neighbors carry label information that an LLM could use. \rag therefore shows these neighbors to the LLM as examples. As with prior \irc\ methods~\cite{colavito2024impact, colavito2025benchmarking}, \rag's most frequent error is labeling questions as bugs, and we target this error by removing the bug-labeled examples for suspected questions (\frag).

In this study, we answer four research questions.

\myparagraph{RQ1} \emph{How effectively can retrieval alone classify issue reports, and at what point does adding more retrieved neighbors stop helping?} The \knn baseline peaks at 59.5\% macro $F_1$ with 15 neighbors from the target project and at 60.4\% with 16 neighbors from all eleven projects. It is weakest on questions and labels 32\% of them as bugs.

\myparagraph{RQ2} \emph{How well does retrieval-augmented few-shot prompting classify issue reports, and how does performance vary with the number of in-context examples?} \rag\ shows the LLM $k$ examples, for seven values of $k$ from 0 (zero-shot prompting) up to 15, the bound set by \knn's peaks. Every configuration with at least one example outperforms both zero-shot prompting and \knn's best result. At each model's best $k$, \rag\ reaches 69.7--76.7\% macro $F_1$ with the target project's labeled issues. Retrieved examples lower the share of questions predicted as bugs from 46--59\% at zero-shot to 26--36\% at the best $k$, but question remains the weakest class.

\myparagraph{RQ3} \emph{Can a balanced retrieval intervention reduce question-to-bug misclassification and improve retrieval-augmented few-shot performance?} The models label many questions as bugs even without examples, so we hypothesize that bug-labeled neighbors retrieved for a question reinforce this error. \Frag\ therefore removes the bug-labeled examples from the prompt, unless they outnumber the question-labeled ones by more than a small margin. It scores higher than \rag\ at every $k{\ge}6$ for all four models, and by 1.5--2.4 points at each method's best $k$ (95\% confidence intervals exclude zero). It also lowers the share of questions predicted as bugs to 19--27\%, at the cost of bug recall.

\myparagraph{RQ4} \emph{How do the two RAG variants compare with LoRA fine-tuning in classification performance, peak GPU memory, and labeled-data needs?} With the same 300 labeled issues of the target project, \rag\ and \frag\ at their best $k$ score higher than LoRA fine-tuning at every model size, by 2.1--6.0 points in macro~$F_1$. When fine-tuning instead uses all 3{,}300 labeled issues of the eleven projects, \rag\ trails it by 2.8 points on average over the four model sizes, and \frag\ trails it by 1.0 point, a difference that is not statistically significant. Once a \knn\ fallback replaces invalid LLM outputs, \frag\ matches fine-tuning (0.1 point ahead on average). Beyond macro~$F_1$, fine-tuning finds more of the actual bugs but over-predicts the bug label, whereas \frag\ makes fewer false bug predictions at most model sizes. The question-to-bug confusion persists under every method we test, including fine-tuning. Both RAG variants need no training, 300 instead of 3{,}300 labeled issues, and 33\% less peak GPU memory than fine-tuning on average.

Overall, retrieval-augmented few-shot prompting is competitive with LoRA fine-tuning while requiring no training, and the two approaches trade different strengths. Which one suits a project depends on its labeled data, its GPU budget, and whether a missed bug or a false bug label costs it more. Our main contributions are the following:

\begin{itemize}
  \item A similarity-weighted $k$-nearest-neighbor baseline (\knn) that measures how much label information the retrieved neighbors carry, bounds the number of in-context examples, and serves as a fallback for invalid LLM outputs.

  \item An evaluation of retrieval-augmented few-shot prompting (\rag) for \irc\ across four model sizes and seven values of $k$.

  \item A training-free retrieval intervention (\frag) that reduces the question-to-bug misclassification known from prior \irc\ work, at the cost of bug recall.

  \item A comparison of \rag and \frag with LoRA fine-tuning of the same models on the same labeled issues, in classification performance, peak GPU memory, and labeled-data needs, reported with paired bootstrap confidence intervals and per-class precision and recall.
\end{itemize}

The rest of the paper is organized as follows. \Cref{sec:03_approach} describes our methods and the fine-tuning baseline, \Cref{sec:setup} the experimental setup, and \Cref{sec:evaluations} the results. \Cref{sec:discussion} analyzes failure cases and practical implications, \Cref{sec:02_related} reviews related work, \Cref{sec:threats} discusses threats to validity, and \Cref{sec:conclusion} concludes and outlines future work.
